\BourbakiExerciseGroup

\begin{enumerate}
\def\labelenumi{\arabic{enumi})}
\tightlist
\item
  \begin{enumerate}
  \def\labelenumii{\alph{enumii})}
  \tightlist
  \item
    设 K 是一个域，f 是 K{[}X{]} 中的一个可分多项式，且 N 是 f 的一个分裂域。证明对于群 \(\operatorname{Gal}(N/K)\) 而言，要使它在 f 于 N 中的根集上可迁地作用，当且仅当 f 在 K{[}X{]} 中不可约.
  \end{enumerate}
\end{enumerate}

\begin{enumerate}
\def\labelenumi{\alph{enumi})}
\setcounter{enumi}{1}
\tightlist
\item
  进一步假设 \(f\) 在 \(\mathbf{K}[\mathbf{X}]\) 中不可约；设 \(x\) 是 \(f\) 在 \(\mathbf{N}\) 中的一个根。要使得不存在域 \(\mathbf{E}\) 满足 \(\mathbf{K} \subset \mathbf{E} \subset \mathbf{K}(x)\) 且异于 \(\mathbf{K}\) 和 \(\mathbf{K}(x)\) ，其充分必要条件是 \(\operatorname{Gal}(\mathbf{N} / \mathbf{K})\) 作为 \(f\) 在 \(\mathbf{N}\) 中的根集上的传递置换群应当是本原的（I，第 142 页，习题 13）。
\end{enumerate}

\begin{enumerate}
\def\labelenumi{\arabic{enumi})}
\setcounter{enumi}{1}
\tightlist
\item
  设 K 为一个域， \(f_{0}\) 为 K{[}X{]} 中一个次数为 n 的不可约且可分的多项式；设 \(\alpha_{i}\) ( \(1 \leqslant i \leqslant n\) ) 为它在 N 中的根，其中 N 是 \(f_{0}\) 的一个分裂域。设 F 为有理分式域 \(\mathrm{N}(\mathbf{X}_{1}, \mathbf{X}_{2}, \ldots, \mathbf{X}_{n})\) ，E 为 F 的子域 \(\mathrm{K}(\mathbf{X}_{1}, \mathbf{X}_{2}, \ldots, \mathbf{X}_{n})\) ；F 是 E 的伽罗瓦扩张，且 \(\operatorname{Gal}(\mathrm{E}/\mathrm{F})\) 典范同构于 \(\operatorname{Gal}(\mathrm{N}/\mathrm{K})\) (V，第 71 页，定理 5；V，第 154 页，习题 8)。证明若 \(\theta = \alpha_{1}\mathbf{X}_{1} + \alpha_{2}\mathbf{X}_{2} + \cdots + \alpha_{n}\mathbf{X}_{n}\) ，则 \(\mathrm{F} = \mathrm{E}(\theta)\) ，并且极小多项式 \(g_{1}\) （ \(\theta\) 在 E{[}T{]} 中的）是该多项式的一个不可约因子
\end{enumerate}

\[
f (T) = \prod_ {\sigma \in \mathfrak {S} _ {n}} \left(T - \alpha_ {1} X _ {\sigma (1)} - \alpha_ {2} X _ {\sigma (2)} - \dots - \alpha_ {n} X _ {\sigma (n)}\right).
\]

每个置换 \(\sigma \in \mathfrak{S}_n\) 定义了 E 的一个 K-自同构，仍记为 \(\sigma\) ，由条件 \(\sigma(X_i) = X_{\sigma(i)}\) （对 \(1 \leqslant i \leqslant n\) ）给定；证明 \(\operatorname{Gal}(N/K)\) 同构于 \(\mathfrak{S}_n\) 的由置换 \(\sigma\) （满足 \(\sigma(g_1) = g_1\) ）组成的子群。进一步地，若 \(f = g_1g_2 \ldots g_r\) 是 \(f\) 在 E{[}T{]} 中的不可约因子分解，证明对每个指标 \(j\) （满足 \(1 \leqslant j \leqslant r\) ）存在一个置换 \(\sigma_j \in \mathfrak{S}_n\) 使得 \(\sigma_j(g_1) = g_j\) .

\begin{enumerate}
\def\labelenumi{\arabic{enumi})}
\setcounter{enumi}{2}
\item
  设 N 是域 K 的无限次伽罗瓦扩张。证明伽罗瓦群 \(\operatorname{Gal}(N/K)\) 的基数大于 N 的子集所成集合的基数（利用 V，第 61 和 62 页）。由此推出存在 \(\operatorname{Gal}(N/K)\) 的非闭子群.
\item
  设 \(\mathbf{N}\) 是域 \(\mathbf{K}\) 的一个伽罗瓦扩张, \((\mathrm{E}_i)_{i\in I}\) 是 \(\mathbf{N}\) 的一族在 \(\mathbf{K}\) 上伽罗瓦的子扩张，使得： \(1^{\circ}\) 对每个指标 \(i\) ，如果 \(\mathbf{F}_i\) 是由 \(\mathbf{E}_j\) （ \(j\neq i\) ）生成的域，则 \(\mathbf{E}_i\cap \mathbf{F}_i = \mathbf{K}\) ; \(2^{\circ}\) \(\mathbf{N}\) 是由诸 \(\mathbf{E}_i\) 的并生成的.
\end{enumerate}

\begin{enumerate}
\def\labelenumi{\alph{enumi})}
\tightlist
\item
  证明： \(\mathbf{N}\) 同构于张量积 \(\otimes_{i\in I}\mathrm{E}_i\) （III，第470页）（定义一个
\end{enumerate}

这个张量积到N上的同构，利用V中第71页的定理5。

\begin{enumerate}
\def\labelenumi{\alph{enumi})}
\setcounter{enumi}{1}
\tightlist
\item
  证明： \(\operatorname{Gal}(\mathbf{N}/\mathbf{K})\) 同构于拓扑群的乘积 \(\operatorname{Gal}(\operatorname{E}_{i}/\operatorname{K})\) .
\end{enumerate}

\begin{enumerate}
\def\labelenumi{\arabic{enumi})}
\setcounter{enumi}{4}
\item
  证明每一个紧致完全不连通群都同构于某个域的伽罗瓦扩张的群Gal(N/K)（利用《一般拓扑学》III，第325页，习题3，化归到该群为有限群的乘积的情形，并同时使用习题4以及V，第59页，例5）。
\item
  设 N 是域 K 的一个伽罗瓦扩张，次数无限，且设 \(\rho, \sigma\) 是 \(\operatorname{Gal}(N/K)\) 的两个不同元素。如果 J 是由 \(x \in N\) （满足 \(\rho(x) \neq \sigma(x)\) ）组成的集合，证明 \(\mathrm{K}(\mathrm{J})\) 是 K 的无限次扩张。（注意，一个域不能是其两个子域的并集，除非它等于其中之一。）
\item
  设 \(\mathbf{N}\) 是域 \(\mathbf{K}\) 的一个伽罗瓦扩张，使得 \(\operatorname{Gal}(\mathbf{N} / \mathbf{K})\) 同构于对称群 \(\mathfrak{S}_n\) (V，第 59 页，例 5)，其中 \(n\) 为非素数的整数；把 \(\operatorname{Gal}(\mathbf{N} / \mathbf{K})\) 与 \(\mathfrak{S}_n\) 等同起来。设 \(\Gamma_1\) 为 \(\mathfrak{S}_n\) 的保持数字 1 不变的、阶为 \((n - 1)!\) 的子群，并设 \(\Gamma_2\) 为 \(\mathfrak{S}_n\) 的阶为 \(n\) 、由循环 \((1, 2, \ldots, n)\) 生成的循环子群（I，第 142 页，习题 12）。设 \(E_1, E_2\) 分别为子群 \(\Gamma_1, \Gamma_2\) 的不变量域。证明 \(E_1\) 与 \(E_2\) 在 \(\mathbf{K}\) 上线性不相交，不存在域 \(F\) （除 \(\mathbf{K}\) 和 \(E_1\) 之外）使得 \(K \subset F \subset E_1\) ，但存在域 \(F'\) （异于 \(E_2\) 和 \(N\) ）使得 \(E_2 \subset F' \subset N\) （利用 I，第 142 页习题 13 及 I，第 143 页习题 14）。
\item
  设 K 是一个域，N 是 K 的有限 n 次伽罗瓦扩张。
\end{enumerate}

\begin{enumerate}
\def\labelenumi{\alph{enumi})}
\item
  设 \(p\) 为一个整除 \(n\) 的素数，并令 \(p^r\) 为 \(p\) 的整除 \(n\) 的最高幂。证明存在子扩张 \(L_i\) （属于 \(N\) ， \(0 \leqslant i \leqslant r\) ）使得对于 \(0 \leqslant i \leqslant r - 1\) , \(L_i\) 包含 \(L_{i+1}\) , \(L_0 = N\) , \(L_i\) 是次数为 \(p\) 的 \(L_{i+1}\) 的伽罗瓦扩张，且 \([L_r: K]\) 不能被 \(p\) 整除（参见 I，第 77 页，定理 1 及 I，第 78 页，定理 2）。
\item
  设 \(x \in \mathbf{N}\) 使得 \(\mathbf{N}\) 由 \(x\) 及其共轭元生成。证明存在一个子扩张 \(\mathbf{E}\) （属于 \(\mathbf{N}\) ）使得 \(x \notin \mathbf{E}\) 且 \([\mathbf{N} : \mathbf{E}] = p'\) 。（考虑 Sylow \(p\) -群 \(\mathrm{H}_j\) （属于 \(\operatorname{Gal}(\mathbf{N} / \mathbf{K})\) ）以及由它们的并生成的（正规）子群 \(\mathbf{H}\) ；若 \(F_j\) （相应地 \(\mathbf{F}\) ）是 \(\mathrm{H}_j\) （相应地 \(\mathbf{H}\) ) 的不变子域，证明 \(x \notin \mathbf{F}\) 并由此推出存在一个指标 \(j\) 使得 \(x \notin F_j\) ；然后我们可以取 \(E = F_{j}\) 。) 得出结论：存在一个子扩张 \(L\) （属于 \(\mathbf{N}\) ）使得 \(L(x)\) 是次数为 \(p\) 的 \(L\) 的伽罗瓦扩张（利用 I，第 77 页，命题 12）。
\item
  反过来，设 \(\Omega\) 为 \(N\) 的一个代数闭包，并假设存在一个子扩张 \(L\) （属于 \(\Omega\) ）使得 \(L(x)\) 是次数为 \(p\) 的、关于 \(L\) 的伽罗瓦扩张。证明此时 \(p\) 整除 \(n\) 。（观察到 \(N \cap L(x) = (N \cap L)(x)\) 且 \((N \cap L)(x)\) 是次数为 \(p\) 的 \(N \cap L\) 的伽罗瓦扩张.)
\end{enumerate}

\begin{enumerate}
\def\labelenumi{\arabic{enumi})}
\setcounter{enumi}{8}
\tightlist
\item
  设 K 是一个特征为 \(p\) 的域, \(\Omega\) 为 K 的一个代数闭扩张， \(x\) 为 \(\Omega\) 的一个元素，使得 \(x \notin K\) ，且 M 为 \(\Omega\) 的不包含 \(x\) 的子扩张所成集合的一个极大元（V，第 148 页，习题 9），因此 \(\Omega\) 是 M 的代数扩张（同上引）。a) 证明：若 M 不是完全域， \(\Omega\) 是 M 的一个 \(p\) -根式扩张，即诸域 \(M(x^{p^{-n}})\) （ \(n \geqslant 0\) ）的并（注意对于每个元素 \(y \in \Omega\) （在 M 上为 \(p\) -根式的），域 \(M(y)\) 包含 \(M(x)\) ）；由此推出我们必然有 \([M(x):M] = p\) ，因为 \(\Omega\) 的每个在 M 上可分的元素必然属于 M）。
\end{enumerate}

\begin{enumerate}
\def\labelenumi{\alph{enumi})}
\setcounter{enumi}{1}
\tightlist
\item
  证明若 M 是完全域， \(\mathbf{M}(x)\) 是 M 的伽罗瓦扩张，次数为素数 q，并且对于每个子扩张 N （ \(\Omega\) 的、在 M 上为有限次伽罗瓦的），其次数 \([N:M]\) 是 q 的幂。（取 q 为整除 \([\mathbf{M}(x):\mathbf{M}]\) 的一个素数；注意 M 的每个不同于 M 且包含于 \(\Omega\) 的扩张都包含 \(\mathbf{M}(x)\) ，并利用习题 8 a)。）
\end{enumerate}

\begin{enumerate}
\def\labelenumi{\arabic{enumi})}
\setcounter{enumi}{9}
\tightlist
\item
  设 \(\mathbf{K}\) 是一个完全域，使得对于每个整数 \(n > 1\) 存在一个伽罗瓦扩张 \(\mathbf{N}\) （ \(\mathbf{K}\) 的）使得 \(\operatorname{Gal}(\mathbf{N} / \mathbf{K})\) 同构于交错群 \(\mathfrak{A}_n\) ；例如，对于每个特征为 0 的域 \(k\) ，域 \(\mathbf{K} = k(\mathbf{X}_n)_{n \in \mathbb{N}}\) （无穷多个未定元上的有理分式域）具有此性质（V，第 59 页，例 5）。
\end{enumerate}

\begin{enumerate}
\def\labelenumi{\alph{enumi})}
\item
  设 A 为一个由 \textgreater{} 1 的整数组成的有限集，且设 \(\Omega\) 为 K 的一个代数闭包。一个元素 \(x \in \Omega\) 称为 A-可构造的，如果存在 K 的扩张组成的一个递增序列 \(K = E_{0} \subset E_{1} \subset \cdots \subset E_{r} \subset \Omega\) ，使得每个次数 \([E_{j}: E_{j-1}]\) （对于 \(1 \leqslant j \leqslant r\) ）都属于 A，且 \(x \in E_{r}\) 。证明存在一个极大元 L ，于子扩张 \(E \subset \Omega\) （其中每个 \(x \in E\) 是 A-可构造的）所成的集合中。对于每个 \(x \notin L\) 证明次数 \([L(x): L]\) 属于 N - A。
\item
  设 \(a\) 是 \(\mathbf{A}\) 的最大元，并且令 \(n\) 为一个满足 \(n \geqslant \max (a + 1,5)\) 的整数。证明存在一个扩张 \(\mathbf{M}\) （ \(\mathbf{L}\) 的）使得 \([\mathbf{M}:\mathbf{L}] = n\) 。（首先观察到，存在一个不可约多项式 \(f \in \mathbf{K}[\mathbf{K}]\) ，其次数为 \(n\) ，使得如果 \(\mathbf{N}\) 是 \(f\) 的分裂域， \(\operatorname{Gal}(\mathbf{N} / \mathbf{K})\) 同构于 \(\mathfrak{A}_n\) ；设 \(y \in \Omega\) 为 \(f\) 在 \(\Omega\) 中的一个根并且令 \(E = K(y)\) 。证明我们必然有 \(E \cap L = K\) ，因此 \(L(E)\) 是所要求的域。用反证法论证，注意到 \(E \cap L\) 的元素是 A-可构造的，并且关系 \(E \cap L \neq K\) 将蕴含一个正规子群 \(\Gamma\) （ \(\mathfrak{A}_n\) 的，满足 \(1 < (\mathfrak{A}_n:\Gamma) < n\) ）的存在性，这将与 \(\mathfrak{A}_n\) 的单性相矛盾（I，第 143 页，习题 16）。
\end{enumerate}

\begin{enumerate}
\def\labelenumi{\arabic{enumi})}
\setcounter{enumi}{10}
\tightlist
\item
  \begin{enumerate}
  \def\labelenumii{\alph{enumii})}
  \tightlist
  \item
    设 E 是域 K 的有限次可分扩张，N 是由 E 生成的 K 的伽罗瓦扩张。若 \(\alpha\) 是 N 的一个元素，其共轭元构成 N 在 K 上的正规基，且 \(\beta = \operatorname{Tr}_{\mathbb{N}/\mathbb{E}}(\alpha)\) ，证明 \(\operatorname{E} = \operatorname{K}(\beta)\) .
  \end{enumerate}
\end{enumerate}

\begin{enumerate}
\def\labelenumi{\alph{enumi})}
\setcounter{enumi}{1}
\tightlist
\item
  设 N 是 K 的有限次伽罗瓦扩张，且令 \(\alpha\) 为 N 的一个元素，其共轭元在 K 上构成一个正规基。若 L 是 N 的一个在 K 上伽罗瓦的子扩张，且 \(\beta = \mathrm{Tr}_{\mathbb{N}/\mathbb{L}}(\alpha)\) ，证明 \(\beta\) 的共轭元构成 L 在 K 上的一个正规基。
\item
  设 L、M 是 K 的两个有限次伽罗瓦扩张，使得 \(L \cap M = K\) ；设 \(\alpha\) （相应地 \(\beta\) ）是 L（相应地，M）中的一个元素，其在 K 上的共轭元构成 L（相应地，M）在 K 上的一个正规基；证明 \(\alpha\beta\) 的在 K 上的共轭元构成 \(\mathrm{K}(\mathrm{L} \cup \mathrm{M})\) 的一个正规基.
\end{enumerate}

\begin{enumerate}
\def\labelenumi{\arabic{enumi})}
\setcounter{enumi}{11}
\tightlist
\item
  设 K 为无限域，N 为 K 的 n 次伽罗瓦扩张， \(f \in K[X]\) 为 N 中一个元素 \(\theta\) 的极小多项式，使得 \(\mathbf{N} = \mathbf{K}(\theta)\) 。群 \(\Gamma = \operatorname{Gal}(N/K)\) 可以作用在 N{[}X{]} 中的多项式上，方式是作用于这些多项式的系数。
\end{enumerate}

\begin{enumerate}
\def\labelenumi{\alph{enumi})}
\tightlist
\item
  设 \(g \in \mathbb{N}[\mathbf{X}]\) 为次数 \(\leqslant n - 1\) 的、使得 \(g(\theta) = 1\) , \(g(\sigma(\theta)) = 0\) （对所有 \(\sigma \neq 1\) ， \(\Gamma\) 中的）的多项式。对于每个子群 \(\Delta\) （ \(\Gamma\) 的），令 \(p_{\Delta} \in \mathbb{N}[\mathbf{X}]\) 为多项式 \(\det(\sigma\tau(g))_{\sigma \in \Delta, \tau \in \Delta}\) 。证明
\end{enumerate}

\[
p _ {\Delta} ^ {2} \equiv \sum_ {\sigma \in \Delta} \sigma (g) \mod f
\]

并推导出 \(p_{\Delta} \neq 0\) .

\begin{enumerate}
\def\labelenumi{\alph{enumi})}
\setcounter{enumi}{1}
\tightlist
\item
  推导出存在一个元素 \(\alpha \in K\) ，使得对于每个子群 \(\Delta\) （ \(\Gamma\) 的）， \(g(\alpha)\) 在子域 \(k(\Delta)\) 上的共轭元构成 N 在 \(k(\Delta)\) 上的一个正规基.
\end{enumerate}

\begin{enumerate}
\def\labelenumi{\arabic{enumi})}
\setcounter{enumi}{12}
\item
  设 E 是域 K 的一个代数扩张， \(E_{s}\) 为 K 的包含在 E 中的最大可分扩张， \(E_{r}\) 为 K 的包含在 E 中的最大 p-根式扩张。在 K 的一个代数闭包中，设 N 是由 E 生成的 K 的拟伽罗瓦扩张；则包含在 N 中的 K 的最大可分扩张是拟伽罗瓦扩张 \(N_{s}\) （由 \(E_{s}\) 生成的）。要使 \(E_{r}\) 是包含在 N 中的 K 的最大 p-根式扩张，必要且充分条件是 E 在 \(E_{r}\) 上可分（参见 V，第 152 页，习题 3）。
\item
  设 \(\mathbf{E}\) 是域 \(\mathbf{K}\) 的一个代数扩张, \(\Gamma\) 为 \(\mathbf{K}\) -自同构组成的群（ \(\mathbf{E}\) 的），且 \(\mathbf{S} \subset \mathbf{E}\) 为 \(\Gamma\) 的不变量域.
\end{enumerate}

\begin{enumerate}
\def\labelenumi{\alph{enumi})}
\item
  证明 E 在 K 上拟伽罗瓦的充分必要条件是 S 是 K 的 p-根式扩张。
\item
  设 \(S_{s}\) 为包含在 S 中的 K 的最大可分扩张。证明 \(S_{s}\) 是满足以下条件的域 F 中的最小者： \(K \subset F \subset E\) 且 E 是 F 的拟伽罗瓦扩张。
\item
  设 \(E_{s}\) 为包含在 E 中的 K 的最大可分扩张；证明 \(\mathrm{E} = \mathrm{S}(\mathrm{E}_{s})\) （注意，E 的 K-自同构中除恒等映射外，没有任何一个会保持 \(\mathrm{S}(\mathrm{E}_{s})\) 的所有元素不变）。
\end{enumerate}

\begin{enumerate}
\def\labelenumi{\arabic{enumi})}
\setcounter{enumi}{14}
\item
  设 E 是域 K 的有限次 n 的伽罗瓦扩张，其伽罗瓦群 G 是幂零的。设 P 为整除 n 的素数集合。则存在 E 的伽罗瓦子扩张 \(E_{p}\) （ \(p \in P\) ）使得 \(\operatorname{Gal}(E_{p}/K)\) 对每个 \(p \in P\) 都是 p-群，并且使得典范同态 \(\otimes E_{p} \to E\) 是一个同构。
\item
  设 \(a\) 是一个整数，并令 \(x_1, x_2, x_3, x_4\) 是 \(P(X) = X^4 - aX - 1\) 在 \(\mathbf{Q}\) 的一个代数闭包中的根。我们把伽罗瓦群 \(G\) （ \(P\) 的）与 \(\{x_1, \ldots, x_4\}\) 上的一个置换群等同起来。假设 \(Q(x_1)\) 不包含二次子扩张（参见 \(V\) ，第 154 页，习题 1）。
\end{enumerate}

\begin{enumerate}
\def\labelenumi{\alph{enumi})}
\item
  证明： \(G = S_{4}\) 或 \(A_{4}\) （利用 V 章第 154 页习题 1：P 不可约，因此 G 传递作用；进一步 \(\mathbf{Q}(x_{1})\) 不包含二次子扩张，因此 \(x_{1}\) 的稳定化子 H 不包含在 G 的指数为 2 的子群中，而 \([G:H] = 4\) ；由此推出 G 不是 2-群，因此其阶可被 3 整除。）
\item
  设 \(\delta = \prod_{1\leqslant i < j\leqslant 4}(x_i - x_j)\) ，且 \(d = \delta^2\) 。证明 \(d = 4^4 -3^3 a^4\) ；进一步有 \(s(\delta) = \varepsilon_s\delta\) ，对
\end{enumerate}

所有 \(s \in G\) 成立。由此推出 \(G = \mathfrak{A}_4\) 如果 \(\delta \in Q\) ，而 \(G = \mathfrak{S}_4\) 否则。证明 \(\delta\) 不是有理的。（因此我们构造了 \(Q\) 上的、具有伽罗瓦群 \(\mathfrak{S}_4\) 的 4 次方程的一个无穷族。）

* 17) 设 \(f\) 是 \(\mathbf{Q}[\mathbf{X}]\) 中的素数次数 \(p\) 的不可约多项式，在 \(\mathbf{C} - \mathbf{R}\) 中恰好有两个根。则伽罗瓦群 \(G\) （ \(f\) 的）同构于 \(\mathfrak{S}_p\) 。（G 包含一个阶为 \(p\) 的元素和一个对换。）*

\begin{enumerate}
\def\labelenumi{\arabic{enumi})}
\setcounter{enumi}{17}
\tightlist
\item
  设 K 为一个域，F 为有理分式域 \(\mathrm{K}(\mathrm{T})\) 。群 \(\mathrm{GL}_{2}(\mathrm{K})\) 按法则 \((\gamma f)(\mathrm{T}) = f\left(\frac{a\mathrm{T} + b}{c\mathrm{T} + d}\right)\) 作用于 F，其中 \(\gamma = \begin{pmatrix} a & b \\ c & d \end{pmatrix} \in \mathrm{GL}_{2}(\mathrm{K})\) 且 \(f \in F\) 。中心 \(K^{*}\) .I （即 \(\mathrm{GL}_{2}(\mathrm{K})\) 的单位元）平凡地作用，因此通过取商，我们得到 \(\mathrm{PGL}_{2}(\mathrm{K}) = \mathrm{GL}_{2}(\mathrm{K}) / \mathrm{K}^{*}\) .I 到 K 的扩张 F 的自同构群内的一个同态。这个同态是双射的（V，第 148 页，习题 11）。
\end{enumerate}

\begin{enumerate}
\def\labelenumi{\alph{enumi})}
\tightlist
\item
  设 H 是 \(\mathrm{PGL}_{2}(\mathbf{K})\) 的一个 h 阶有限子群，并设 E 为 H 的不变量域，从而 \(F = E(T)\) 是 E 的伽罗瓦扩张，其伽罗瓦群为 H（V，第 66 页，定理 3）。设
\end{enumerate}

\[
\mathrm{H} (\mathrm{X}) = \prod_ {s \in \mathrm{H}} (\mathrm{X} - s (\mathrm{T})) = \sum_ {i = 0} ^ {h} (- 1) ^ {i} \mathrm{S} _ {i} (\mathrm{T}) \mathrm{X} ^ {h - i}
\]

为 T 在 E 上的极小多项式。对于每个满足 \(0 \leqslant i \leqslant h\) 的 i，或者 \(S_{i}(T) \in K\) ，或者 \(E = K(S_{i}(T))\) 。（利用 V，第 157 页，习题 11 a)。）

\begin{enumerate}
\def\labelenumi{\alph{enumi})}
\setcounter{enumi}{1}
\item
  对于每个 \(s \in H\) ，令 \(\begin{pmatrix} a_{s} & b_{s} \\ c_{s} & d_{s} \end{pmatrix}\) 为它在 \(GL_{2}(K)\) 中的一个代表元，并且令 \(H_{0}\) 为 H 中由满足 \(c_{s} = 0\) 的那些 s 组成的子群；设 \(h_{0} = CardH_{0}\) 并令 \(R^{H}(T) = S_{h_{0}}(T)\) 。证明 \(E = K(R^{H}(T))\) 且 \(R^{H}(T) = P(T)/Q(T)^{h_{0}}\) ，其中 P 是次数为 h 的多项式，且 \(Q(T) = \prod_{s} (c_{s}T + d_{s})\) ，其中 s 遍历陪集 \(H_{0}s\) （异于 \(H_{0}\) 的）的一个代表系；因此 Q 的次数为 \((h/h_{0}) - 1\) .
\item
  假设 \(K^{*}\) 包含一个有限阶 n 的子群 \(\mu_{n}\) 。设 \(C_{n}\) 为（在 \(PGL_{2}(K)\) 中的）\(GL_{2}(K)\) 的由矩阵 \(\begin{pmatrix} a & 0 \\ 0 & 1 \end{pmatrix}, a \in \mu_{n}\) 组成的子群的像。则 \(C_{n}(X) = X^{n} - T^{n}\) 且 \(R^{C_{n}}(T) = (-1)^{n+1}T^{n}\) .
\item
  设 \(w\) 为 \(\mathrm{PGL}_2(\mathbf{K})\) 中由矩阵 \(\begin{pmatrix} 0 & 1 \\ 1 & 0 \end{pmatrix}\) 表示的元素。则 \(w^2 = 1\) ，且 \(wsw^{-1} = s^{-1}\) （若 \(s\) 由一个对角矩阵表示）。子群 \(\mathrm{D}_n\) （ \(\mathrm{PGL}_2(\mathbf{K})\) 的、由 \(\mathrm{C}_n\) 与 \(w\) 生成的）因此是 \(2n\) 阶的。我们有 \(\mathrm{D}_n(\mathbf{X}) = (\mathbf{X}^n - \mathrm{T}^n)\left(\mathbf{X}^n - \left(\frac{1}{\mathrm{T}}\right)^n\right) = \mathbf{X}^{2n} + (-1)\mathbf{R}^{\mathrm{D}_n(\mathrm{T})} + 1\) ，其中 \(\mathbf{R}^{\mathrm{D}_n(\mathrm{T})} = -((-T)^n + (-T)^{-n})\) .
\item
  设 S 为 \(\mathrm{PGL}_2(\mathbf{K})\) 中由矩阵 \(\left( \begin{array}{cc}0 & -1\\ 1 & 1 \end{array} \right)\) 的像生成的 3 阶群。我们有 \(S(X) = (X - T)\left(X - \frac{1}{1 - T}\right)\left(X - \frac{T - 1}{T}\right)\) 且 \(R^s (T) = \frac{T^3 - 3T + 1}{T^2 - 1}\) .
\item
  设 K 是一个含有 q 个元素的有限域。令 U 为（在 \(\mathrm{PGL}_{2}(\mathbf{K})\) 中的）由所有矩阵 \(\begin{pmatrix}1 & 0 \\ b & 1\end{pmatrix}\) , \(b \in K\) 构成的群的像。则
\end{enumerate}

\[
\mathrm{U} (\mathrm{X}) = (\mathrm{X} - \mathrm{T}) ^ {q} - (\mathrm{X} - \mathrm{T}) = \mathrm{X} ^ {q} - \mathrm{X} - (\mathrm{T} ^ {q} - \mathrm{T}) \quad \text { and } \quad \mathrm{R} ^ {\mathrm{U}} (\mathrm{T}) = \mathrm{T} ^ {q} - \mathrm{T}.
\]

设 B 是由 \(C_{q-1}\) 与 U 生成的群；其阶为 \(q(q-1)\) 。我们有

\[
\mathrm{B} (\mathrm{X}) = \left(\mathrm{X} ^ {q} - \mathrm{X}\right) ^ {q - 1} - \left(\mathrm{T} ^ {q} - \mathrm{T}\right) ^ {q - 1} \quad \text { and } \quad \mathrm{R} ^ {\mathrm{H}} (\mathrm{T}) = \left(\mathrm{T} ^ {q} - \mathrm{T}\right) ^ {q - 1}.
\]

\begin{enumerate}
\def\labelenumi{\arabic{enumi})}
\setcounter{enumi}{18}
\item
  设 E 是域 K 的伽罗瓦扩张，伽罗瓦群为 G。对每个 \(x \in E\) 用 Gx 表示 x 在 E 中的共轭元素的集合；这是一个有限的 G-集。设 x，y 为 E 中的两个元素。要使 G-集 Gx 和 Gy 同构（即存在一个双射 \(Gx \rightarrow Gy\) ，与 G 在这两个集合上的作用相容），其必要且充分条件是扩张 \(\mathrm{K}(x)\) 与 \(\mathrm{K}(y)\) 共轭。
\item
  设 E 是域 K 的一个代数扩张，使得 K{[}X{]} 中每个非常数多项式在 E 中至少有一个根。证明 E 是代数闭的。（在 E 的代数闭包中——从而也是在 K 的代数闭包中——设 F 是 K 的有限次扩张。要证明 \(F \subset E\) ，可归结为 F 是拟伽罗瓦扩张的情形，然后利用 V 第 76 页命题 13，再归结为 F 是伽罗瓦扩张或 p-根式扩张的情形。）
\item
  设 \(f(\mathbf{X}) = \mathbf{X}^{n} + a_{1}\mathbf{X}^{n-1} + \cdots + a_{n}\) 是 Z{[}X{]} 中的多项式，且 p 为素数。假设 \(a_{i} \equiv 0 \pmod{p}\) , \(1 \leqslant i \leqslant n\) ，且 \(a_{n} \not\equiv 0 \pmod{p^{2}}\) 。证明 \(f(\mathbf{X})\) 在 Q{[}X{]} 中不可约。（设 \(f(\mathbf{X}) = \mathrm{P}(\mathbf{X}) \mathrm{Q}(\mathbf{X})\) 是 Q{[}X{]} 中的一个分解。证明：若 P 和 Q 为首一多项式，则它们具有整数系数。（模 p 约化。）
\end{enumerate}

* 22) 设 \(m\) 是一个偶数 \(>0\) , \(r\) 为一个整数 \(\geqslant 3\) ，且 \(n_1 < \cdots < n_{r-2}\) 为 \(r-2\) 个偶数。令 \(g(\mathbf{X}) = (\mathbf{X}^2 + m)(\mathbf{X} - n_1) \ldots (\mathbf{X} - n_{r-2})\) .

\begin{enumerate}
\def\labelenumi{\alph{enumi})}
\item
  令 \(f(\mathbf{X}) = g(\mathbf{X}) - 2\) 。证明 \(f(\mathbf{X})\) 至少有 \(r - 2\) 个实根。计算 \(f\) 在 \(\mathbf{C}\) 中的根的平方和。由此推出，对于 \(m \geqslant \sum n_i^2 / 2\) , \(f(\mathbf{X})\) 是一个具有整数系数的多项式，它有 \(r - 2\) 个实根和两个非实的共轭复根。
\item
  证明： \(f(\mathbf{X})\) 在 \(\mathbf{Q}[\mathbf{X}]\) 中不可约。（写 \(f(\mathbf{X}) = \mathbf{X}' + \sum_{k=1}^{r} a_k \mathbf{X}'^{-k}\) ；证明
\end{enumerate}

\(a_{k}\) 是偶数且 \(a_{r}\) 不能被4整除；应用习题21。）

\begin{enumerate}
\def\labelenumi{\alph{enumi})}
\setcounter{enumi}{2}
\item
  证明当 \(r\) 为素数且 \(m \geqslant \sum n_i^2 / 2\) 时，分裂域的伽罗瓦群 \(G\) （ \(f\) 的）同构于 \(\mathfrak{S}_r\) （参见 V，第 158 页，习题 17）。
\item
  证明当 \(r = 4\) 且 \(m \geqslant \sum n_i^2 / 2\) 时，分裂域在 \(\mathbf{Q}\) 上的次数等于 8 或 24，并且它等于 8 当且仅当 \(n_1 = -n_2\) ，或者也当且仅当 \(f(\mathbf{X})\) 具有形式 \(\mathrm{P}(\mathbf{X}^2)\) ，其中 \(\mathbf{P}\) 是一个次数为 2 的多项式。*
\end{enumerate}

\begin{enumerate}
\def\labelenumi{\arabic{enumi})}
\setcounter{enumi}{22}
\tightlist
\item
  设 \(k\) 是一个特征为2的域。
\end{enumerate}

\begin{enumerate}
\def\labelenumi{\alph{enumi})}
\tightlist
\item
  设 \(A_{1}, \ldots, A_{n}\) 为未定元，K 为域 \(k(\mathbf{A}_{1}, \ldots, \mathbf{A}_{n})\) ，E 为多项式 \(P = X^{n} + A_{1}X^{n-1} + \cdots + A_{n} \in K[X]\) 的一个分裂域，且 \(X_{1}, \ldots, X_{n}\) 为 P 在 E 中的根，因此 \(E = k(X_{1}, \ldots, X_{n})\) （V，第 21 页）。我们将 E/K 的伽罗瓦群与对称群 \(S_{n}\) （通过变量 \(X_{1}, \ldots, X_{n}\) 的置换而作用）等同起来（V，第 59 页）。设
\end{enumerate}

\[
\beta \left(\mathrm{X} _ {1}, \dots , \mathrm{X} _ {n}\right) = \sum_ {i <   j} \mathrm{X} _ {i} / \left(\mathrm{X} _ {i} + \mathrm{X} _ {j}\right) \in \mathrm{E} \text { and } \mathrm{C} \left(\mathrm{A} _ {1}, \dots , \mathrm{A} _ {n}\right) = \sum_ {i <   j} \left(\mathrm{X} _ {i} \mathrm{X} _ {j}\right) / \left(\mathrm{X} _ {i} ^ {2} + \mathrm{X} _ {j} ^ {2}\right) \in \mathrm{K}.
\]

则 \(\beta \notin K\) 且

\[
\boldsymbol {\beta} ^ {2} + \boldsymbol {\beta} + \mathbf {C} = 0.
\]

我们有 \(\sigma(\beta) = \beta\) （对 \(\sigma \in \mathfrak{A}_n\) ），且 \(\sigma(\beta) = \beta + 1\) （对 \(\sigma \in \mathfrak{S}_n - \mathfrak{A}_n\) ）。由此推出 \(K(\beta)\) 是 \(E/K\) 的唯一的二次子扩张.

\begin{enumerate}
\def\labelenumi{\alph{enumi})}
\setcounter{enumi}{1}
\tightlist
\item
  设 \(D \in Z[T_{1}, \ldots, T_{n}]\) 为多项式 \(X^{n} + T_{1}X^{n-1} + \cdots + T_{n}\) 的判别式。证明存在 Q, R， \(S \in Z[T_{1}, \ldots, T_{n}]\) ，满足
\end{enumerate}

\[
\mathbf {D} = \mathbf {Q} ^ {2} - 4 \mathbf {R} + 8 \mathbf {S}
\]

且

\[
\mathrm{C} = \mathrm{C} (\mathrm{A} _ {1}, \dots , \mathrm{A} _ {n}) = \frac {\mathrm{R} (\mathrm{A} _ {1} , \dots , \mathrm{A} _ {n})}{\mathrm{Q} ^ {2} (\mathrm{A} _ {1} , \dots , \mathrm{A} _ {n})} = \frac {\mathrm{R} (\mathrm{A} _ {1} , \dots , \mathrm{A} _ {n})}{\mathrm{D} (\mathrm{A} _ {1} , \dots , \mathrm{A} _ {n})}.
\]

对于 \(\mathbf{Z}[T_1, \ldots, T_n]\) 的元素 U, V, W 的每一个满足 \(D = U^2 - 4V + 8W\) 的三元组，存在 \(M \in \mathbf{Z}[A_1, \ldots, A_n]\) 使得

\[
C = \frac {V (A _ {1} , \dots , A _ {n})}{D (A _ {1} , \dots , A _ {n})} + \left(\frac {M}{D (A _ {1} , \dots , A _ {n})}\right) ^ {2} + \frac {M}{D (A _ {1} , \dots , A _ {n})}.
\]

\begin{enumerate}
\def\labelenumi{\alph{enumi})}
\setcounter{enumi}{2}
\tightlist
\item
  设 \(f = X^{n} + a_{1}X^{n-1} + \cdots + a_{n}\) 是 k{[}X{]} 的一个可分多项式，且 G 是一个分裂扩张的伽罗瓦群。证明 \((a_{1}, \ldots, a_{n})\) 在 C 中是可代入的，并且以下条件等价
\end{enumerate}

\begin{enumerate}
\def\labelenumi{(\roman{enumi})}
\item
  \(G\) 的每个元素诱导 \(f\) 的根的一个偶置换;
\item
  存在 \(\alpha \in k\) 使得 \(C(a_{1},\ldots ,a_{n}) = \alpha^{2} + \alpha\) .
\end{enumerate}

* \(d\) ) 今后假设 \(k\) 是有限的，并任意选取 \(\mathbf{Z}[T_1, \ldots, T_n]\) 的元素 U、V、W 以满足 \(b\) ) 中所述的条件.

证明 \(c\) ) 的 (i) 和 (ii) 与以下条件等价：

\begin{enumerate}
\def\labelenumi{(\roman{enumi})}
\setcounter{enumi}{2}
\tightlist
\item
  \[
\operatorname{Tr} _ {k / \mathbf {F} _ {2}} (\mathbf {C}) = 0;
\]
\item
  \[
\mathrm{Tr} _ {k / \mathbf {F} _ {2}} \left(\frac {\mathbf {V} (a _ {1} , \dots , a _ {n})}{\mathbf {D} (a _ {1} , \dots , a _ {n})}\right) = 0.
\]
\end{enumerate}

\begin{enumerate}
\def\labelenumi{(\alph{enumi})}
\setcounter{enumi}{21}
\tightlist
\item
  \(f\) 的不可约因子的个数与 \(n\) 模 2 同余（使用 V，第 165 页，习题 22）。*
\end{enumerate}

